% Unified notation for Bernoulli types - Latent/Observed Framework
% ===================================================================
% This file emphasizes the fundamental distinction between latent (true)
% values and their observations throughout the textbook.

% CORE PRINCIPLE: Latent values need no decoration; observed values use tilde

% ========== OBSERVATION NOTATION ==========

% The universal observation operator
\newcommand{\obs}[1]{\widetilde{#1}}

% Explicit markers (for clarity in definitions)
\newcommand{\latent}[1]{#1}
\newcommand{\observed}[1]{\widetilde{#1}}

% ========== TYPE NOTATION ==========

% Basic types
\newcommand{\Type}[1]{\mathtt{#1}}
\newcommand{\Bool}{\Type{Bool}}
\newcommand{\Nat}{\mathbb{N}}
\newcommand{\Real}{\mathbb{R}}
\newcommand{\Unit}{\Type{Unit}}
\newcommand{\Void}{\Type{Void}}
\newcommand{\Maybe}[1]{\Type{Maybe}[#1]}

% Boolean values
\newcommand{\True}{\mathtt{true}}
\newcommand{\False}{\mathtt{false}}

% ========== BERNOULLI TYPES ==========

% Bernoulli type constructor: observations of type T with order n error model
\newcommand{\bernoulli}[2]{\mathcal{B}\langle #1, #2 \rangle}
\newcommand{\Bernoulli}[2]{\mathcal{B}\langle #1, #2 \rangle}

% Common Bernoulli types
\newcommand{\BernoulliBool}{\bernoulli{\Bool}{}}
\newcommand{\BernoulliSet}[1]{\bernoulli{\mathcal{P}(#1)}{}}
\newcommand{\BernoulliMap}[2]{\bernoulli{#1 \to #2}{}}

% ========== SET NOTATION ==========

% Sets: no decoration for latent sets
% Example: S is a latent set, \obs{S} is an observed/approximate set

% Power set
\newcommand{\PS}[1]{\mathcal{P}(#1)}
\newcommand{\PowerSet}[1]{\mathcal{P}(#1)}
\newcommand{\EmptySet}{\emptyset}

% Set operations (standard LaTeX)
% Use \cup, \cap, \setminus directly
\newcommand{\SetComplement}[1]{\overline{#1}}
\newcommand{\Complement}[1]{\overline{#1}}
\newcommand{\SetIndicator}[1]{\mathbf{1}_{#1}}
\newcommand{\Indicator}[1]{\mathbf{1}_{#1}}

% ========== FUNCTION NOTATION ==========

% Functions: no decoration for latent functions
% Example: f is a latent function, \obs{f} is an observed/approximate function

% Function symbols
% \newcommand{\mapsto}{\rightarrow} % Already defined in LaTeX
\newcommand{\pfun}{\rightharpoonup}  % Partial function
\newcommand{\compose}{\circ}

% Function properties
\newcommand{\Dom}{\mathsf{dom}}
\newcommand{\Cod}{\mathsf{cod}}
\newcommand{\Range}{\mathsf{range}}
\newcommand{\Image}{\mathsf{img}}

% ========== PROBABILITY NOTATION ==========

% Basic probability
\newcommand{\Prob}[1]{\mathbb{P}\left[#1\right]}
\newcommand{\ProbCond}[2]{\mathbb{P}\left[#1 \mid #2\right]}
\newcommand{\Expect}[1]{\mathbb{E}\left[#1\right]}
\newcommand{\ExpectCond}[2]{\mathbb{E}\left[#1 \mid #2\right]}
\newcommand{\Var}[1]{\mathrm{Var}\left[#1\right]}
\newcommand{\Cov}[2]{\mathrm{Cov}\left[#1, #2\right]}

% Channel notation: observed | latent
\newcommand{\Channel}[2]{\ProbCond{\text{observe } #1}{\text{latent } #2}}

% ========== ERROR RATES ==========

\newcommand{\fprate}{\alpha}      % False positive rate
\newcommand{\fnrate}{\beta}       % False negative rate
\newcommand{\error}{\epsilon}     % Generic error rate
\newcommand{\errorrate}{\epsilon}

% Performance metrics
\newcommand{\FPR}{\mathsf{FPR}}
\newcommand{\FNR}{\mathsf{FNR}}
\newcommand{\TPR}{\mathsf{TPR}}
\newcommand{\TNR}{\mathsf{TNR}}
\newcommand{\PPV}{\mathsf{PPV}}  % Positive predictive value
\newcommand{\NPV}{\mathsf{NPV}}  % Negative predictive value

% ========== INFORMATION THEORY ==========

\newcommand{\Entropy}[1]{H(#1)}
\newcommand{\MutualInfo}[2]{I(#1; #2)}
\newcommand{\KL}[2]{D_{\mathrm{KL}}(#1 \| #2)}

% ========== COMPLEXITY ==========

\newcommand{\BigO}[1]{O(#1)}
\newcommand{\BigOmega}[1]{\Omega(#1)}
\newcommand{\BigTheta}[1]{\Theta(#1)}
\newcommand{\littleo}[1]{o(#1)}
\newcommand{\littleomega}[1]{\omega(#1)}
\newcommand{\softO}[1]{\tilde{O}(#1)}

% ========== COMMON FUNCTIONS ==========

\newcommand{\Card}[1]{\lvert#1\rvert}  % Cardinality
\newcommand{\Floor}[1]{\lfloor #1 \rfloor}
\newcommand{\Ceil}[1]{\lceil #1 \rceil}
\newcommand{\Norm}[1]{\|#1\|}

% ========== LOGIC OPERATIONS ==========

% Use standard LaTeX: \land, \lor, \neg
\newcommand{\AND}{\wedge}
\newcommand{\OR}{\vee}
\newcommand{\NOT}{\neg}
\newcommand{\XOR}{\oplus}
\newcommand{\IMPLIES}{\Rightarrow}
\newcommand{\IFF}{\Leftrightarrow}

% ========== SPECIAL SYMBOLS ==========

\newcommand{\Given}{\mid}
\newcommand{\Ind}{\perp\!\!\!\perp}  % Independence

% ========== CODE FORMATTING ==========

\newcommand{\code}[1]{\texttt{#1}}
\newcommand{\codecomment}[1]{\textit{\small // #1}}

% ========== ALGORITHM NOTATION ==========

\newcommand{\Assignment}{\leftarrow}
\newcommand{\Random}{\stackrel{\$}{\leftarrow}}
% \newcommand{\Repeat}{\text{repeat}} % Already defined by algorithm2e
% \newcommand{\Until}{\text{until}} % Already defined by algorithm2e
\newcommand{\Return}{\text{return}}

% ========== COMMON ABBREVIATIONS ==========

\newcommand{\iid}{\text{i.i.d.}}
\newcommand{\wrt}{\text{w.r.t.}}
\newcommand{\st}{\text{s.t.}}
\newcommand{\ie}{\textit{i.e.}}
\newcommand{\eg}{\textit{e.g.}}
\newcommand{\cf}{\textit{cf.}}
\newcommand{\viz}{\textit{viz.}}

% ========== MATRIX NOTATION ==========

\newcommand{\Matrix}[1]{\mathbf{#1}}
\newcommand{\Vector}[1]{\mathbf{#1}}
\newcommand{\Transpose}{^\top}

% ========== DEPRECATED NOTATION ==========
% These are kept for backward compatibility but should be phased out:

\newcommand{\Set}[1]{#1}  % Just use S instead of \Set{S}
\newcommand{\ASet}[1]{\obs{#1}}  % Use \obs{S} instead
\newcommand{\Fun}[1]{\mathsf{#1}}  % Just use f
\newcommand{\AFun}[1]{\obs{\mathsf{#1}}}  % Use \obs{f}

% ========== COMPARISON OPERATORS ==========

\newcommand{\equals}{==}
\newcommand{\notequals}{!=}
\newcommand{\assign}{=}